\song{Sáro! (Traband)}
%\re{ } \textit{Refrén po každé sloce}
\cho 
\ch{Ami}Sáro, \ch{Emi}Sáro, \ch{F}v noci se mi \ch{C}zdálo, \\
že \ch{F}tři andělé \ch{C}Boží k nám \ch{F}přišli na o\ch{G}běd.\\ 
\ch{Ami}Sáro, \ch{Emi}Sáro, jak \ch{F}moc a nebo \ch{C}málo \\
mi \ch{F}chybí abych \ch{C}tvojí duši \ch{F}mohl rozu\ch{G}m{ě}t.

\ve
Sbor kajícných mnichů jde krajinou v tichu \\
a pro všechnu lidskou pýchu má jen přezíravý smích \\
Z prohraných válek se vojska domů vrací \\
ač zbraně stále burácí a bitva zuří v nich.
\cho
\ve
Vévoda v zámku čeká na balkóně \\
až přivedou mu koně, pak mává na pozdrav \\
Srdcová dáma má v každé ruce růže, \\
tak snadno pohřbít může sto urozených hlav.
\cho
\ve
Královnin šašek s pusou od povidel \\
sbírá zbytky jídel a myslí na útěk \\
A v podzemí skrytí slepí alchymisté \\
už objevili jistě proti povinnosti lék

\cho
Sáro, Sáro, jak moc a nebo málo \\
ti chybí abys mojí duši mohla rozumět 

\ve
Páv pod tvým oknem zpívá sotva procit \\
o tajemstvích noci ve tvých zahradách \\
A já, potulný kejklíř, co svázali mu ruce, \\
teď hraju o tvé srdce a chci mít tě nadosah!
\cho
Sáro, Sáro, pomalu a líně \\
s hlavou na tvém klíně chci se probouzet \\
Sáro, Sáro, Sáro, Sáro rosa padá ráno \\
a v poledne už možná bude jiný svět\\

Sáro, Sáro, vstávej milá Sáro \\
andělé k nám \ch{Dmi}přišli na \ch{C}oběd

\esong